\section{Interaction-resolved Quantum Optimal Control}
\label{sec:results}

\noindent
\noindent
We now use the interaction coordinates above to formulate interaction-resolved quantum optimal control objectives and apply them to a realistically parametrized solid-state spin register based on a nitrogen-vacancy (NV) center in diamond. For a given control pulse, the system is evolved under the effective Hamiltonian to obtain the propagator, from which computational-basis phases are extracted in the corresponding diagonalizing frame and transformed into interaction coordinates used as optimization targets.


\medskip
\noindent
We present two three-qubit targets that test the same interaction-resolved control principle. We first synthesize a diagonal tripartite entangler by targeting the desired three-body coordinate while suppressing all pairwise coordinates and leaving local coordinates unconstrained. We then synthesize a non-diagonal three-qubit entangler by performing the optimization in a suitable diagonalizing frame, demonstrating that the same scheme applies to interactions generated by non-diagonal Pauli strings.

\subsection{Control Objectives in Interaction-Resolved Coordinates}
\label{sec:control_objective_phase_invariants}

\noindent
Target interaction structures can be specified by assigning desired values $\{\phi_S^\star\}$ to selected interaction coordinates associated with the relevant many-body terms, as represented in Eq.~\eqref{eq:phi_expansion_u_diag}.

% \medskip
% \noindent
% The decomposition involves interaction contributions supported on all subsets $S \subseteq \{1,\dots,n\}$ and thus comprises $2^n$ independent coefficients, reflecting the full parameter count of a general $n$-qubit unitary in the diagonal frame up to a global phase. In many relevant settings, however, only a restricted set of $k$-body interaction terms determines the desired operation. In such cases, the relevant parameter space is reduced to a subset of support-selective phase invariants associated with the corresponding interaction structure. Specifying a target operation therefore amounts to fixing a set of invariants $\{\phi_S^\star\}$ associated with selected subsets $S \subseteq \{1,\ldots,n\}$ corresponding to $k$-body contributions on certain qubits.\\

\noindent
To quantify deviations from desired many-body interaction contributions, a distance measure can be defined directly on the selected interaction coordinates rather than on the full unitary.\\
Given target values for a chosen set of qubit subsets  $\mathcal{S}_{\mathrm{target}}$, we define the interaction-resolved control cost function
\begin{equation}
\label{eq:phase-invariant-cost}
\mathcal{J}
=
\sum_{S\in\mathcal{S}_{\mathrm{target}}}
w_S
\Bigl[
1-\cos\!\bigl(2(\phi(S)-\phi_S^\star)\bigr)
\Bigr],
\end{equation}
where $w_S$ determines the relative importance of the corresponding interaction terms. 
Furthermore, this formulation also enables interaction-selective optimization: only the interaction contributions associated with subsets in $\mathcal{S}_{\mathrm{target}}$ enter the objective function. For example, leaving degrees of freedom, such as local terms, unconstrained, enlarges the space of admissible solutions, allowing the optimization to exploit this flexibility to find more efficient implementations.

\medskip
\noindent
\textbf{Remark.}
The factor of two in Eq.~\eqref{eq:phase-invariant-cost} reflects the $\pi$-periodicity of Pauli-$Z$ interaction phases. Indeed,
\begin{equation}
e^{-i(\theta+\pi)Z_S}
=
-\,e^{-i\theta Z_S},
\label{eq:pi_peridicity}
\end{equation}
so shifting the interaction phase by $\pi$ changes only the collective sign of the operator without modifying its physical action. The objective therefore identifies interaction phases differing by integer multiples of $\pi$ as physically equivalent.


\subsection{Model Spin Hamiltonian}
\label{sec:spin_model}

The following Hamiltonian model, describing the dynamics of a spin register based on a nitrogen-vacancy (NV) center in diamond, serves as the working model for the simulated platform, for which we optimize the controls. A full derivation of this Hamiltonian starting from the laboratory-frame NV Hamiltonian is given in Appendix~\ref{app:nv_model}. In this platform an electronic spin is coupled to nearby nuclear spins, forming a controllable multi-qubit system driven by microwave control fields. In the interaction picture with respect to the free precessing parts of the  Hamiltonian and after applying the rotating-wave approximation (RWA) at the electronic resonance, the Hamiltonian takes the form
\begin{equation}
\begin{aligned}
&\notag H_{\mathrm{NV}}(t)\\
&=\frac{\Omega(t)}{2}\sum_{\mathbf m}
\big[\sigma_x\cos((\omega_{\mathrm{MW}}-\Lambda(\mathbf m))\,t)\\
&\hspace{1.0cm}-\sigma_y\sin((\omega_{\mathrm{MW}}-\Lambda(\mathbf m))\,t)\big] \otimes \ket{\mathbf m}\!\bra{\mathbf m} \\
&-\;\frac{\Omega(t)}{2}\big[\sigma_x\sin((\omega_{\mathrm{MW}}-\Lambda_s) t)+\sigma_y\cos((\omega_{\mathrm{MW}}-\Lambda_s) t)\big] \\
&\hspace{1.0cm}\otimes\sum_i \theta_i \big[I'_{ix}\cos(\delta_i t-\phi_i)+I'_{iy}\sin(\delta_i t-\phi_i)\big],
\end{aligned}
\label{eq:H_NV_RWA_main}
\end{equation}
The first term describes the coherent microwave drive of the electronic qubit conditioned on the nuclear configuration, while the second term captures the weak microwave-mediated coupling to nearby nuclear spins.  
Here, $\Omega(t)$ is the Rabi frequency that can be tuned via a microwave magnetic field (with carrier frequency $\omega_{\mathrm{MW}}$) and serves as our control pulse, $\Lambda_s= D - \gamma_e B_0$ describes the bare electron transition frequency between its two ground-state levels \(\ket{0_e}\!\leftrightarrow\!\ket{-1_e}\), including zero-field splitting $D$ and the Zeeman shift due to the static magnetic field $B_0$ with the gyromagnetic ratio $\gamma_e$. $\Lambda(\mathbf{m})$ describes the configuration-dependent transition frequencies of the electron, with  $\mathbf m=(m_N,\{m_{C_i}\})$ labeling each nuclear spin register configuration, i.e.\ a joint eigenstate of the nuclear spin operators along the quantization axis, where the quantum numbers $m_N$ and $m_{C_i}$ correspond to the eigenvalues of the spin operator $I_z$ for the $^{14}$N and $^{13}$C nuclei. 
%The function $E(t)$ denotes the time-dependent control pulse envelope defined in Eq.~\eqref{eq:pulse_parametrization}.
The quantities $\delta_i= \frac{\gamma_i B_0 + \omega_i}{2}$ describe weak effective transverse modulations of the nuclear spin. The parameters $\theta_i$ and $\phi_i$ characterize the local hyperfine misalignment of the nuclear spins.\\
The specific physical parameters used in the numerical simulations of the NV register are given in Appendix~\ref{app:nv_params}.


\subsection{Pulse Parametrization and Optimization}
\label{sec:pulse_parametrization} 
We control the system by optimizing the temporal shape of the Rabi frequency appearing in Eq.~\eqref{eq:H_NV_RWA_main}. We start by the parametrization 
\begin{equation}
\Omega(t) = \chi(t)\sum_{i=1}^{N_c} a_{i} \cos(\omega_{c,i} t + \phi_{c,i}),
\label{eq:pulse_parametrization}
\end{equation}
with frequency amplitudes $a_{i}$ carrying the physical units of the control envelope, control frequencies $\omega_{c,i}$, and control phases $\phi_{c,i}$.
The resulting parameters are collected into a time-independent parameter 
\begin{equation}
\mathbf p =
(a_{1},\ldots,a_{N_c},\;
\omega_{c,1},\ldots,\omega_{c,N_c},\;
\phi_{c,1},\ldots,\phi_{c,N_c}),
\label{eq:pulse_parametrization_example}
\end{equation}
which are to be optimized iteratively. All optimized pulse parameters, resulting in the following, are provided in the public repository~\cite{baran2026_multiqubit_control}.
The multiplicative factor $\chi(t)$ is a Tukey (tapered cosine) window that smoothly turns the pulse on and off,
\begin{equation}
\chi(t)=\tfrac12\left[1-\cos\!\left(2\pi\frac{t}{T}\right)\right], 
\qquad t\in[0,\alpha T],
\end{equation}
The parameter $\alpha\in[0,1]$ controls the taper fraction, we generally set it to $\alpha=0.15$.

\noindent
This analytic parametrization~\eqref{eq:pulse_parametrization}, used, similarly, for example in CRAB~\cite{caneva2011chopped,rach2015dressing}, GRAPE~\cite{khaneja2005optimal,de2011second}, and hybrid approaches~\cite{motzoi2011optimal,sorensen2018quantum,sorensen2020optimization,preti2022continuous}, provides a compact and experimentally compatible representation of the control fields. 


Thus, for a given parameter vector $\mathbf p$, the control pulse
$\Omega(t;\mathbf p)$ defines the effective propagator
\begin{equation}
U_{\mathrm{pulse}}(\mathbf p,T)
=
\mathcal{T}
\exp\!\left(
-i\int_0^T
H_{\mathrm{NV}}(\Omega(t;\mathbf p),t)\,dt
\right).
\label{eq:U_pulse_definition}
\end{equation}





\subsection{Diagonal Three-Qubit Entangler}
\label{sec:results_nv_ccz}
